\setcounter{framenumber}{0}
%25

\begin{frame}
  \titlepage
\end{frame}

\begin{frame}{Learning Goals}
By the end of this lecture, students should be able to:
\begin{itemize}
    \item define the moment generating function (MGF);
    \item compute MGFs of common distributions;
    \item extract moments from MGFs;
    \item understand how MGFs determine distributions;
    \item compute distributions of sums using MGFs;
    \item derive moments of Exponential and Normal distributions using MGFs.
\end{itemize}
\end{frame}

\begin{frame}{Definition of the MGF}
\begin{keybox}{Definition}
The moment generating function (MGF) of $X$ is
\[
{
M_X(t)=\mathbb{E}(e^{tX})
}
\]
provided the expectation is finite in some open interval around $0$.

\vspace{0.2cm}

Equivalently,

\[
M_X(t)
=
\begin{cases}
\displaystyle
\sum_x e^{tx} p_X(x),
& \text{if $X$ is discrete},\\[3ex]

\displaystyle
\int_{-\infty}^{\infty} e^{tx} f_X(x)\,dx,
& \text{if $X$ is continuous}.
\end{cases}
\]
\end{keybox}

\vspace{0.25cm}

\begin{keybox}{Quick check}
Every valid MGF satisfies
\[
M_X(0)=1.
\]
\end{keybox}
\end{frame}






% Convention: q=1-p.
% Geom(p): number of failures before the first success, so k=0,1,2,\dots
% NegBin(r,p): number of failures before the r-th success.

\begin{frame}{MGFs: Discrete Distributions I}
\small
\[
M_X(t)=\mathbb{E}(e^{tX})=\sum_x e^{tx}p_X(x).
\]

\begin{keybox}{Bernoulli}
If \(X\sim\operatorname{Bern}(p)\), then
\[
M_X(t)=e^{t\cdot 1}p+e^{t\cdot 0}(1-p)=pe^t+q.
\]
\end{keybox}

\begin{keybox}{Discrete Uniform}
If \(X\sim\operatorname{DUnif}\{1,\dots,m\}\), then
\[
M_X(t)=\frac1m\sum_{k=1}^m e^{tk}
=\frac1m\left(e^t+e^{2t}+\cdots+e^{mt}\right)
=\frac{e^t(e^{mt}-1)}{m(e^t-1)}.
\]
\end{keybox}

\end{frame}





\begin{frame}{MGFs: Discrete Distributions II}

\small

\begin{keybox}{Binomial}
If \(X\sim\operatorname{Bin}(n,p)\), then
\[
M_X(t)=\sum_{k=0}^n e^{tk}\binom{n}{k}p^kq^{n-k}
=\sum_{k=0}^n \binom{n}{k}(pe^t)^kq^{n-k}
=(q+pe^t)^n.
\]
\end{keybox}

\begin{keybox}{Poisson}
If \(X\sim\operatorname{Pois}(\lambda)\), then
\[
M_X(t)=\sum_{k=0}^{\infty} e^{tk}e^{-\lambda}\frac{\lambda^k}{k!}
=e^{-\lambda}\sum_{k=0}^{\infty}\frac{(\lambda e^t)^k}{k!}
=e^{-\lambda}e^{\lambda e^t}
=e^{\lambda(e^t-1)}.
\]
\end{keybox}

\end{frame}



\begin{frame}{MGFs: Discrete Distributions III}

\small

\begin{keybox}{Geometric}
If \(X\sim\operatorname{Geom}(p)\), where \(X\) counts failures before the first success, then
\[
M_X(t)=\sum_{k=0}^{\infty} e^{tk}q^kp
=p\sum_{k=0}^{\infty}(qe^t)^k
=\frac{p}{1-qe^t},
\qquad qe^t<1.
\]
\end{keybox}

\begin{keybox}{Negative Binomial}
If \(X\sim\operatorname{NegBin}(r,p)\), where \(X\) counts failures before the \(r\)-th success, then
\[
M_X(t)=\sum_{k=0}^{\infty}e^{tk}\binom{k+r-1}{k}p^rq^k
=p^r\sum_{k=0}^{\infty}\binom{k+r-1}{k}(qe^t)^k
=\left(\frac{p}{1-qe^t}\right)^r.
\]
\end{keybox}

\end{frame}



\begin{frame}{MGFs: Discrete Distributions IV}

\small

\begin{keybox}{Hypergeometric}
If
\[
X\sim\operatorname{HGeom}(w,b,n),
\]
then
\[
\Prob(X=k)=
\frac{\binom{w}{k}\binom{b}{n-k}}{\binom{w+b}{n}},
\]
where
\[
\max(0,n-b)\le k\le \min(w,n).
\]

Therefore,
\[
M_X(t)
=
\sum_{k=\max(0,n-b)}^{\min(w,n)}
e^{tk}
\frac{\binom{w}{k}\binom{b}{n-k}}{\binom{w+b}{n}}.
\]
\end{keybox}

\textbf{Note}: Unlike Binomial or Poisson, this MGF usually does not simplify to a nice elementary closed form.

\end{frame}




\begin{frame}{MGFs: Continuous Distributions I}

\small

\[
M_X(t)=\mathbb{E}(e^{tX})=\int_{-\infty}^{\infty} e^{tx}f_X(x)\,dx.
\]

\begin{keybox}{Uniform}
If \(X\sim\operatorname{Unif}(a,b)\), then
\[
M_X(t)=\int_a^b e^{tx}\frac1{b-a}\,dx
=\frac1{b-a}\left[\frac{e^{tx}}{t}\right]_a^b
=\frac{e^{tb}-e^{ta}}{t(b-a)},
\qquad t\ne0.
\]
Also, \(M_X(0)=1\).
\end{keybox}

\end{frame}







\begin{frame}{MGFs: Continuous Distributions II}

\small

\begin{keybox}{Exponential}
If \(X\sim\operatorname{Exp}(\lambda)\), then
\[
M_X(t)=\int_0^\infty e^{tx}\lambda e^{-\lambda x}\,dx
=\lambda\int_0^\infty e^{-(\lambda-t)x}\,dx
=\frac{\lambda}{\lambda-t},
\qquad t<\lambda.
\]
\end{keybox}

\begin{keybox}{Gamma}
If \(X\sim\operatorname{Gamma}(\alpha,\lambda)\), with shape \(\alpha\) and rate \(\lambda\), then
\[
M_X(t)
=
\int_0^\infty e^{tx}
\frac{\lambda^\alpha}{\Gamma(\alpha)}
x^{\alpha-1}e^{-\lambda x}\,dx
=
\frac{\lambda^\alpha}{\Gamma(\alpha)}
\int_0^\infty x^{\alpha-1}e^{-(\lambda-t)x}\,dx.
\]
Let \(u=(\lambda-t)x\). Then \(x=u/(\lambda-t)\), \(dx=du/(\lambda-t)\), and
\[
M_X(t)
=
\frac{\lambda^\alpha}{\Gamma(\alpha)}
\frac{1}{(\lambda-t)^\alpha}
\underbrace{
\int_0^\infty u^{\alpha-1}e^{-u}\,du
}_{\Gamma(\alpha)}
=
\frac{\lambda^\alpha}{(\lambda-t)^\alpha}
=
\left(\frac{\lambda}{\lambda-t}\right)^\alpha,
\qquad t<\lambda.
\]
\end{keybox}

\end{frame}

























\begin{frame}{MGF of $N(0,1)$}

\small

\begin{keybox}{Standard Normal}
Let \(Z\sim N(0,1)\). Then
\[
M_Z(t)=\mathbb{E}(e^{tZ})
=\frac1{\sqrt{2\pi}}\int_{-\infty}^{\infty}e^{tz}e^{-z^2/2}\,dz.
\]
Complete the square:
\[
tz-\frac{z^2}{2}
=
-\frac12(z^2-2tz)
=
-\frac12\bigl((z-t)^2-t^2\bigr)
=
-\frac{(z-t)^2}{2}+\frac{t^2}{2}.
\]
Therefore,
\[
M_Z(t)
=
\frac1{\sqrt{2\pi}}\int_{-\infty}^{\infty}
e^{-(z-t)^2/2}e^{t^2/2}\,dz
=
e^{t^2/2}
\frac1{\sqrt{2\pi}}\int_{-\infty}^{\infty}e^{-(z-t)^2/2}\,dz.
\]
Since the last integral equals \(\sqrt{2\pi}\),
\[
{M_Z(t)=e^{t^2/2}}.
\]
\end{keybox}

\end{frame}


















\begin{frame}{MGF of $N(\mu,\sigma^2)$}

\small

\begin{keybox}{Normal}
Let \(X\sim N(\mu,\sigma^2)\). Write
\[
X=\mu+\sigma Z,
\qquad Z\sim N(0,1).
\]
Then
\[
M_X(t)
=
\mathbb{E}(e^{tX})
=
\mathbb{E}\left(e^{t(\mu+\sigma Z)}\right)
=
e^{\mu t}\mathbb{E}(e^{\sigma t Z}).
\]
Using
\[
M_Z(s)=e^{s^2/2},
\]
with \(s=\sigma t\), we get
\[
M_X(t)
=
e^{\mu t}M_Z(\sigma t)
=
e^{\mu t}e^{\sigma^2t^2/2}
=
{e^{\mu t+\frac12\sigma^2t^2}}.
\]
\end{keybox}

\end{frame}



































\lecturepart{Moments from MGFs}
\begin{frame}{MGFs Generate Moments}
\begin{keybox}{Theorem}
The $n^{th}$ population moment is obtained by differentiating the MGF:
\[
\mathbb{E}(X^n)=M_X^{(n)}(0).
\]
\end{keybox}

\vspace{0.25cm}

Examples:
\[
\mathbb{E}(X)=M_X'(0),
\qquad
\mathbb{E}(X^2)=M_X''(0).
\]
\end{frame}


\begin{frame}{Proof: MGFs Generate Moments}
Recall that the moment generating function is
\[
M_X(t)=\mathbb{E}(e^{tX}).
\]

Differentiate once:
\[
M_X'(t)
=
\frac{d}{dt}\mathbb{E}(e^{tX})
=
\mathbb{E}\left(\frac{d}{dt}e^{tX}\right)
=
\mathbb{E}(Xe^{tX}).
\]

Now plugging in $t=0$,
\[
M_X'(0)
=
\mathbb{E}(Xe^{0})
=
\mathbb{E}(X).
\]

Differentiate another time noting that $X$ is a constant w.r.t. $t$.
\[
M_X''(t)
=
\frac{d}{dt}\mathbb{E}(Xe^{tX})
=
\mathbb{E}\left(\frac{d}{dt}Xe^{tX}\right) = \mathbb{E}(X^2 e^{tX}).
\]
Plugging in $t=0$,
\[
M_X''(0)
=
\mathbb{E}(X^2 e^0)
=
\mathbb{E}(X^2).
\]

Similarly, differentiating $n$ times gives
\[
M_X^{(n)}(t)=\mathbb{E}(X^n e^{tX}).
\]

Therefore, at $t=0$,
\[
M_X^{(n)}(0)
=
\mathbb{E}(X^n e^0)
=
\mathbb{E}(X^n).
\]
\end{frame}



%------------------------------------------------
\begin{frame}{Taylor Expansion Around a Point}

Let \(f\) be infinitely differentiable in a neighborhood of \(a\).
Its \(n\)th-order Taylor polynomial around \(a\) is
\[
T_n(x)
=
\sum_{k=0}^{n}
\frac{f^{(k)}(a)}{k!}(x-a)^k.
\]

Taylor's theorem gives
\[
f(x)
=
\sum_{k=0}^{n}
\frac{f^{(k)}(a)}{k!}(x-a)^k
+
R_n(x),
\]
where, under the usual differentiability assumptions, the remainder can
be written as
\[
R_n(x)
=
\frac{f^{(n+1)}(\xi)}{(n+1)!}(x-a)^{n+1}
\]
for some \(\xi\) between \(a\) and \(x\).

\begin{keybox}{Infinite Taylor Series}
If $R_n(x)\longrightarrow 0
\qquad\text{as }n\to\infty,$ then
\[
f(x)
=
\sum_{k=0}^{\infty}
\frac{f^{(k)}(a)}{k!}(x-a)^k.
\]
\end{keybox}

\end{frame}

%------------------------------------------------
\begin{frame}{Taylor Expansion of the Exponential Function}

For
\[
f(u)=e^u,
\]
we have
\[
f^{(k)}(u)=e^u
\qquad\text{for every }k\geq 0.
\]

Expanding around \(a=0\),
\[
f^{(k)}(0)=1,
\]
so
\[
e^u
=
\sum_{k=0}^{\infty}\frac{u^k}{k!}
=
1+u+\frac{u^2}{2!}+\frac{u^3}{3!}+\cdots.
\]

The exponential function is analytic on all of \(\mathbb{R}\), so this
series converges to \(e^u\) for every \(u\in\mathbb{R}\).

Now substitute
\[
u=tX.
\]
For every fixed realization of \(X\),
\[
e^{tX}
=
\sum_{k=0}^{\infty}\frac{(tX)^k}{k!}
=
1+tX+\frac{t^2X^2}{2!}
+\frac{t^3X^3}{3!}+\cdots.
\]

Thus, the identity holds \textbf{pointwise} for every value of \(X\).

\end{frame}

%------------------------------------------------
\begin{frame}{Taking Expectations: Why Is It Valid?}

The moment-generating function of \(X\) is
\[
M_X(t)=\mathbb{E}\!\left(e^{tX}\right),
\]
provided this expectation is finite.

Using the Taylor expansion,
\[
M_X(t)
=
\mathbb{E}\left[
\sum_{k=0}^{\infty}\frac{t^kX^k}{k!}
\right].
\]

To write
\[
\mathbb{E}\left[
\sum_{k=0}^{\infty}\frac{t^kX^k}{k!}
\right]
=
\sum_{k=0}^{\infty}
\frac{t^k}{k!}\mathbb{E}(X^k),
\]
we must justify interchanging the expectation and the infinite sum.

A sufficient condition is that, for some \(r>|t|\),
\[
\mathbb{E}\!\left(e^{r|X|}\right)<\infty.
\]
\end{frame}

%------------------------------------------------
\begin{frame}{Moments as Taylor Coefficients of the MGF}
 

Therefore,
\[
\boxed{
M_X(t)
=
\sum_{k=0}^{\infty}
\frac{\mathbb{E}(X^k)}{k!}t^k
}
\]
whenever the required interchange is valid.

Equivalently,
\[
M_X(t)
=
1
+t\mathbb{E}(X)
+\frac{t^2}{2!}\mathbb{E}(X^2)
+\frac{t^3}{3!}\mathbb{E}(X^3)
+\cdots.
\]

Comparing this with the Taylor expansion
\[
M_X(t)
=
\sum_{k=0}^{\infty}
\frac{M_X^{(k)}(0)}{k!}t^k,
\]
we obtain
\[
\boxed{
M_X^{(k)}(0)=\mathbb{E}(X^k)
}
\]
whenever \(M_X\) exists in an open neighborhood of \(0\).

\begin{keybox}{Main Idea}
The \(k\)th raw moment of \(X\) is the \(k\)th derivative of its MGF
evaluated at \(t=0\).
\end{keybox}

\end{frame}


















\begin{frame}{Example: Bernoulli Moments}
For
\[
M_X(t)=pe^t+q,
\]
we compute:
\[
M_X'(t)=pe^t.
\]

Thus:
\[
M_X'(0)=p.
\]

So
\[
{
\mathbb{E}(X)=p.
}
\]
\end{frame}
\begin{frame}{MGF Uniqueness Theorem}

\begin{keybox}{Theorem (Uniqueness of MGF)}
Let $X$ and $Y$ be random variables. Suppose there exists some $h>0$ such that
\[
M_X(t)=\mathbb{E}(e^{tX})<\infty,
\qquad
M_Y(t)=\mathbb{E}(e^{tY})<\infty
\]
for every $t\in(-h,h)$.

If
\[
M_X(t)=M_Y(t)
\qquad \text{for all } t\in(-h,h),
\]
then
\[
X\overset{d}{=}Y,
\]
that is, $X$ and $Y$ have the same distribution.
\end{keybox}

\vspace{0.2cm}

\textbf{Why is this useful?} If we want to determine the distribution of a random variable $X$:

\[
\boxed{
\text{Find }M_X(t)
\;\longrightarrow\;
\text{recognize it as the MGF of a known distribution}
\;\longrightarrow\;
\text{identify the distribution of }X.
}
\]
\end{frame}


\begin{frame}{Using the MGF Uniqueness Theorem}

Suppose we compute the MGF of a random variable $X$ and obtain
\[
M_X(t)
=
\exp\left(\mu t+\frac{\sigma^2t^2}{2}\right),
\qquad t\in(-h,h)
\]
for some $h>0$.

\vspace{0.2cm}

We recognize this as the MGF of a normal random variable:
\[
Y\sim\mathcal N(\mu,\sigma^2)
\qquad\Longrightarrow\qquad
M_Y(t)
=
\exp\left(\mu t+\frac{\sigma^2t^2}{2}\right).
\]

Therefore,
\[
M_X(t)=M_Y(t)
\qquad\text{for all }t\in(-h,h).
\] a 

By the \textbf{MGF Uniqueness Theorem},
\[
\boxed{X\overset{d}{=}Y}
\qquad\Longrightarrow\qquad
\boxed{X\sim\mathcal N(\mu,\sigma^2).}
\]

\vspace{0.3cm}

\begin{keybox}{Key Idea}
If the MGF of an unknown random variable matches the MGF of a known
distribution on an interval containing $0$, then the random variable
must have that distribution.
\end{keybox}

\end{frame}



\lecturepart{MGFs of Sums}

\begin{frame}{MGF of a Sum}
\begin{keybox}{Theorem}
If $X$ and $Y$ are independent, then
\[
{
M_{X+Y}(t)
=
M_X(t)\times M_Y(t).
}
\]
More generally, if $X_1,\dots, X_n$ are independent, then
$$
M_{\sum_{i=1}^n X_i}(t)
=
\prod_{i=1}^nM_{X_i}(t)
$$
\end{keybox}

This is one of the main reasons MGFs are useful. Since using the MGF Uniqueness theorem, the MGF characterize the distribution in full.
\end{frame}

\begin{frame}{Binomial MGF}
A Binomial random variable can be written as:
\[
X=I_1+\cdots+I_n,
\]
where the $I_j$ are independent Bernoulli$(p)$ random variables.

\vspace{0.25cm}

Since each Bernoulli MGF is
\[
pe^t+q,
\]
we get
\[
{
M_X(t)
=
(pe^t+q)^n.
}
\]
\end{frame}

\begin{frame}{Negative Binomial MGF}
If
\[
X\sim \operatorname{NBin}(r,p),
\]
then
\[
X
=
G_1+\cdots+G_r,
\]
where the $G_j$ are independent Geometric$(p)$ random variables.

\vspace{0.25cm}

Thus,
\[
{
M_X(t)
=
\left(
\frac{p}{1-qe^t}
\right)^r.
}
\]
\end{frame}



\lecturepart{Sums via MGFs}

\begin{frame}{Sum of Independent Poissons}
If
\[
X\sim \operatorname{Pois}(\lambda),
\qquad
Y\sim \operatorname{Pois}(\mu),
\]
independently, then
\[
M_X(t)=e^{\lambda(e^t-1)},
\]
and
\[
M_Y(t)=e^{\mu(e^t-1)}.
\]

Therefore,
\[
M_{X+Y}(t)
=
e^{(\lambda+\mu)(e^t-1)}.
\]

Hence:
\[
{
X+Y\sim \operatorname{Pois}(\lambda+\mu).
}
\]

\textbf{Remark}: This can be generalized to $n$ independent Poisson r.v.s using induction.
\end{frame}

\begin{frame}{Sum of Independent Normals}
If
\[
X_1\sim N(\mu_1,\sigma_1^2),
\qquad
X_2\sim N(\mu_2,\sigma_2^2),
\]
independently, then
\[
M_{X_1+X_2}(t)
=
e^{(\mu_1+\mu_2)t
+\frac12(\sigma_1^2+\sigma_2^2)t^2}.
\]

Thus,
\[
{
X_1+X_2
\sim
N(\mu_1+\mu_2,\sigma_1^2+\sigma_2^2).
}
\]


\textbf{Remark}: This can be generalized to $n$ independent Normal r.v.s using induction.

\end{frame}







\begin{frame}{Distributions Closed Under Addition \textbf{Under Independence}}

\begin{center}
\small
\renewcommand{\arraystretch}{1.45}
\begin{tabular}{lll}
\toprule
\textbf{Distribution} & \textbf{Assumption} & \textbf{Sum Distribution} \\
\midrule

Bernoulli &
$X_i\sim \operatorname{Bern}(p)$ &
$\displaystyle \sum_{i=1}^n X_i
\sim \operatorname{Bin}(n,p)$
\\

Binomial &
$X_i\sim \operatorname{Bin}(n_i,p)$ &
$\displaystyle \sum_{i=1}^n X_i
\sim
\operatorname{Bin}\!\left(\sum n_i,p\right)$
\\

Poisson &
$X_i\sim \operatorname{Pois}(\lambda_i)$ &
$\displaystyle \sum_{i=1}^n X_i
\sim
\operatorname{Pois}\!\left(\sum \lambda_i\right)$
\\

Negative Binomial &
$X_i\sim \operatorname{NegBin}(r_i,p)$ &
$\displaystyle \sum_{i=1}^n X_i
\sim
\operatorname{NegBin}\!\left(\sum r_i,p\right)$
\\

Geometric &
$X_i\sim \operatorname{Geom}(p)$ &
$\displaystyle \sum_{i=1}^n X_i
\sim
\operatorname{NegBin}(n,p)$
\\

Normal &
$X_i\sim N(\mu_i,\sigma_i^2)$ &
$\displaystyle \sum_{i=1}^n X_i
\sim
N\!\left(\sum \mu_i,\sum \sigma_i^2\right)$
\\

Gamma &
$X_i\sim \operatorname{Gamma}(\alpha_i,\lambda)$ &
$\displaystyle \sum_{i=1}^n X_i
\sim
\operatorname{Gamma}\!\left(\sum \alpha_i,\lambda\right)$
\\

Exponential &
$X_i\sim \operatorname{Exp}(\lambda)$ &
$\displaystyle \sum_{i=1}^nX_i
\sim
\operatorname{Gamma}(n,\lambda)$
\\

\bottomrule
\end{tabular}
\end{center}

\vspace{0.2cm}

\small
All random variables above are assumed independent.
Some families additionally require a common parameter.

\end{frame}



\begin{frame}{Summary}
\begin{itemize}
    \item MGFs encode moments through derivatives.

    \item MGFs determine distributions.

    \item Independent sums correspond to products of MGFs.

    \item MGFs simplify moment calculations dramatically.

    \item The Normal and Exponential MGFs are especially important.
\end{itemize}

\vspace{0.25cm}

\begin{keybox}{Next lecture}
Joint distributions.
\end{keybox}
\end{frame}